% SGA 3, Expose V, Loop 1 English reconstruction.
% Scope: Lemma 6.1 proof (a), proof (b), and opening of proof (c).
% Authority: Polo--Gille Exp5-13oct24.pdf, component-local p. 22
% (combined-reader p. 284), re-edited page header 272.
% Jacob Reinhold's CC BY 4.0 Markdown is comparison/drafting material only.

另一方面, 若 \(d_0:X_1\to X_0\) 开 (分别地, 普遍闭),
则由 \(d_0\) 作基变换得到的 \(u\) 也有同样性质.
由定理~\hyperref[thm:sga3-V-4.1]{4.1}, \(q\) 满射、整且开.
可知 \(r\) 满射, 并且当 \(d_0\) 开 (分别地, 普遍闭) 时也开 (分别地, 普遍闭).
由于 \(v\) 满射, \(p\) 也有同样性质.
这就证明了\hyperref[lem:sga3-V-6.1-ip]{(i') 项}.

现在假设 \(S\) 局部 Noether, \(X_0\) 在 \(S\) 上局部有限型,
故 \(X_0\) 局部 Noether.

证明 \(Y\) 在 \(S\) 上局部有限呈示. 设
\[
  S'=\operatorname{Spec}R
\]
是 \(S\) 的仿射开集, 设
\[
  Y'=\operatorname{Spec}B
\]
是 \(Y\) 中映入 \(S'\) 的仿射开集, 并设
\[
  U'=\operatorname{Spec}A
\]
是 \(Y'\) 在 \(U\) 中的逆像. 由于 \(R\) 是 Noether 环,
只须证明 \(B\) 是有限生成的 \(R\)-代数.
由 \hyperref[sec:sga3-V-4]{\S\S4} 和 \hyperref[sec:sga3-V-5]{5},
\(B\) 包含在 \(A\) 中, 而后者是有限生成的 \(R\)-代数;
因此, 由于 \(R\) 是 Noether 环, 且 \(A\) 在 \(B\) 上整, 得到所述断言.

最后, 由于 \(X_0\to S\) 局部有限型, \(p\) 也局部有限型
(EGA~I, 6.6.6). 由于 \(Y\) 局部 Noether, 故 \(p\) 局部有限呈示.

还须证明\hyperref[lem:sga3-V-6.1-ii]{(ii) 项}的最后一个断言.
再假设 \(X_0\) 在 \(S\) 上有限型.
由于 \(p\) 满射, \(Y\) 在 \(S\) 上拟紧, 因而在 \(S\) 上有限型.
由于 \(S\) 局部 Noether, 态射 \(X_0\to S\) 和 \(Y\to S\)
都有限呈示, 因而 \(p:X_0\to Y\) 也有限呈示
(EGA~IV\(_1\), 1.6.2(v)).

\medskip
\phantomsection\label{par:sga3-V-6-b}%
\noindent\textbf{b) \hyperref[lem:sga3-V-6.1-iii]{(iii)} 的证明.}---%
以 \(V\) 为底的群胚 \(V_\ast\), 既同构于 \(U_\ast\)
在基变换 \(u\) 下的逆像, 也同构于 \(X_\ast\)
在基变换 \(v\) 下的逆像. 因此有一个双重笛卡儿方块:
\SGAThreeDiagram[0.65\textwidth]
  {figures/exp5/p272_lemma61_double_cartesian_square.png}
  {\hyperref[lem:sga3-V-6.1-iii]{引理~6.1(iii)} 的证明中所用的双重笛卡儿方块.}
  {fig:sga3-V-6.1-double-cartesian-square}
  {第 V 篇, 分篇第~22 页 / 合订本第~284 页 / 重版本页眉页码~272}
由于 \(u_0\boxtimes u_1\) 满射, \(v_0\boxtimes v_1\) 也满射.
由于 \(v\times v\) 满射, 复合态射
\[
  V_1\longrightarrow X_0\times_Y X_0
\]
满射, 因而 \(d_0\boxtimes d_1\) 也满射.

\medskip
\phantomsection\label{par:sga3-V-6-c}%
\noindent\textbf{c) \hyperref[lem:sga3-V-6.1-i]{(i)} 的证明.}---%
还须证明 \((Y,p)\) 是所有赋环空间范畴中 \((d_0,d_1)\) 的余核.
先证明 \(Y\) 由 \(X_0\) 通过等同如下的点 \(x\) 和 \(y\) 得到:
存在 \(z\in X_1\), 使得
\[
  d_0(z)=x,\qquad d_1(z)=y.
\]
事实上, \(p\) 满射, 且 \(p d_0=p d_1\).
反过来, 若 \(p(x)=p(y)\), 则存在点 \(z'\in X_0\times_Y X_0\),
其第一投影为 \(x\), 第二投影为 \(y\).
若 \(z\in X_1\) 满足
\[
  (d_0\boxtimes d_1)(z)=z',
\]
则 \(d_0(z)=x\) 且 \(d_1(z)=y\), 正如所求.

另一方面, 若 \(W\) 是 \(X_0\) 的饱和开集,
则 \(W\cap U\) 是 \(U\) 的饱和开集.
由定理~\hyperref[thm:sga3-V-4.1]{4.1}, \(q(W\cap U)\) 在 \(Y\) 中开.
由于 \(q(W\cap U)=p(W)\), \(Y\) 上的拓扑是由 \(X_0\) 诱导的商拓扑.

还须证明环层的典范序列
\begin{equation}
  \mathcal O_Y
  \longrightarrow p_\ast\mathcal O_{X_0}
  \rightrightarrows
  p_\ast d_{0\ast}\mathcal O_{X_1}
  =
  p_\ast d_{1\ast}\mathcal O_{X_1}
  \label{eq:sga3-V-6.1-sheaf-exactness}
\end{equation}
正合.

% Translation coverage: entire assigned component; adjoining components retain the original statement/proof seams.
% Source lines: 103; translated body lines before this audit comment: 96.
% Math preserved: 78 inline formulas, 6 display blocks, 1 equation environments; mathematical prose translated where present.
% Labels: 3/3; hyperrefs: 9/9; SGARefLinks: 0/0; inline targets: 0/0; footnotes: 0/0.
% Figure arguments preserved; all rendered captions and natural-language source locators translated.
% Unresolved TODOs: 0; compilation and visual inspection remain with the volume integrator.
